\documentclass[12pt,a4paper]{article}
\usepackage{goyabase}
\usepackage[utf8]{inputenc}
\usepackage[T1]{fontenc}
\usepackage[spanish, es-noshorthands, es-tabla]{babel}
% xcolor ya lo carga goyabase, no lo declaramos de nuevo para evitar Option clash
\usepackage{tabularx}
\usepackage{enumitem}
\usepackage{listings}
\usepackage{csquotes}

% Definicion de colores que podrian faltar
\definecolor{grisclaro}{gray}{0.9}

% Configuracion de listings para comandos SQL y Bash
\lstset{
    basicstyle=\ttfamily\small,
    breaklines=true,
    frame=single,
    backgroundcolor=\color{gray!10},
    keywordstyle=\color{blue},
    commentstyle=\color{green!40!black},
    stringstyle=\color{orange},
    showstringspaces=false,
    inputencoding=utf8,
    extendedchars=true,
    literate={á}{{\'a}}1 {é}{{\'e}}1 {í}{{\'i}}1 {ó}{{\'o}}1 {ú}{{\'u}}1 {ñ}{{\~n}}1 {¿}{{\textquestiondown}}1 {¡}{{\textexclamdown}}1
}

\begin{document}

\section{Memoria técnica}

\subsection{Comandos a utilizar:}

\subsubsection{Crear copia de seguridad completa}
\begin{lstlisting}[language=bash]
mysqldump -u admin -p1234 logistica_global > 00_copia_inicial.sql
\end{lstlisting}

\subsubsection{Cargar la copia de seguridad completa}
\begin{lstlisting}[language=bash]
mysql -u admin -p1234 logistica_global < 00_copia_inicial.sql
\end{lstlisting}

\subsubsection{Crear una copia de seguridad incremental}
Rotamos los logs con:
\begin{lstlisting}[language=bash]
mysql -u admin -p1234 -e "FLUSH LOGS;" 
\end{lstlisting}

Generamos la copia de seguridad:
\begin{lstlisting}[language=bash]
mysqlbinlog --read-from-remote-server -u admin -p1234 -h localhost binlog.??????
\end{lstlisting}

\subsection{Proceso de creación:}

\subsubsection{1.1 Copia de seguridad completa}
\begin{lstlisting}[language=bash]
mysqldump -u admin -p1234 logistica_global > 00_copia_inicial.sql
\end{lstlisting}

\subsubsection{Rotación de logs}
\begin{lstlisting}[language=bash]
mysql -u admin -p1234 -e "FLUSH LOGS;" 
\end{lstlisting}

\subsubsection{Resuelvo el ejercicio 1}
\begin{lstlisting}[language=bash]
mysql -u admin -p1234 logistica_global < limpieza/p1_asir.sql
\end{lstlisting}

\subsubsection{Rotación de logs}
\begin{lstlisting}[language=bash]
mysql -u admin -p1234 -e "FLUSH LOGS;" 
\end{lstlisting}
para garantizar que no hay más ejecuciones sql en la copia de seguridad.

\subsubsection{¿Qué log me interesa?}
\begin{lstlisting}[language=bash]
mysql -u admin -p1234 -e "SHOW MASTER STATUS;" 
\end{lstlisting}
Respuesta: binlog.000948 (en este entorno).

Por tanto en el binlog.000948 tengo la copia incremental del ejercicio 1.

\subsubsection{1.2.a Creo la primera incremental}
\begin{lstlisting}[language=bash]
mysqlbinlog --read-from-remote-server -u admin -p1234 -h localhost binlog.000948 > 01_incremental_ejercicio1.sql
\end{lstlisting}

\subsubsection{Resuelvo el ejercicio 2}
\begin{lstlisting}[language=bash]
mysql -u admin -p1234 logistica_global < limpieza/p2_asir.sql 
\end{lstlisting}

\subsubsection{Rotación de logs}
Cerramos el log 949 antes de hacer la copia.
\begin{lstlisting}[language=bash]
mysql -u admin -p1234 -e "FLUSH LOGS;" 
\end{lstlisting}

\subsubsection{1.2.b Creo la incremental}
\begin{lstlisting}[language=bash]
mysqlbinlog --read-from-remote-server -u admin -p1234 -h localhost binlog.000949 > 02_incremental_ejercicio2.sql
\end{lstlisting}

\subsubsection{1.3 Creo la diferencial}
\begin{lstlisting}[language=bash]
mysqlbinlog --read-from-remote-server -u admin -p1234 -h localhost binlog.000948 binlog.000949 > 03_diferencial_ejers_1_y_2.sql
\end{lstlisting}

\subsubsection{Resuelvo el ejercicio 3}
\begin{lstlisting}[language=bash]
mysql -u admin -p1234 logistica_global < limpieza/p3_asir.sql 
\end{lstlisting}

\subsubsection{Rotación de logs}
Cierro el log 950.
\begin{lstlisting}[language=bash]
mysql -u admin -p1234 -e "FLUSH LOGS;" 
\end{lstlisting}

\subsubsection{1.4 Creo la incremental:}
\begin{lstlisting}[language=bash]
mysqlbinlog --read-from-remote-server -u admin -p1234 -h localhost binlog.000950 > 03_incremental_ejercicio3.sql
\end{lstlisting}

\subsection{1.5 Resumen de copias de seguridad}
\begin{itemize}
    \item \textbf{00\_copia\_inicial.sql}: copia de seguridad completa antes de resolver la limpieza.
    \item \textbf{01\_incremental\_ejercicio1.sql}: copia de seguridad incremental desde la bbdd original que incluye la resolución del ejercicio 1.
    \item \textbf{02\_incremental\_ejercicio2.sql}: copia de seguridad incremental desde la bbdd original que incluye la resolución del ejercicio 2.
    \item \textbf{03\_diferencial\_ejers\_1\_y\_2.sql}: copia de seguridad diferencial que combina las incrementales de los ejercicios 1 y 2.
    \item \textbf{03\_incremental\_ejercicio3.sql}: copia de seguridad incremental desde la bbdd original que incluye la resolución del ejercicio 3.
\end{itemize}

\subsection{Restaurar las copias de seguridad}
Situado desde la carpeta donde están los backups:

\begin{lstlisting}[language=bash]
mysql -u admin -p1234 logistica_global < 00_copia_inicial.sql
\end{lstlisting}

Comprobamos que no hay la columna de coste:
\begin{lstlisting}[language=bash]
mysql -u admin -p1234 logistica_global -e "select coste_eur from mantenimientos_flota"
\end{lstlisting}

Restauramos la primera incremental:
\begin{lstlisting}[language=bash]
mysql -u admin -p1234 logistica_global < 01_incremental_ejercicio1.sql
\end{lstlisting}

Comprobamos que sí hay "coste\_eur":
\begin{lstlisting}[language=bash]
mysql -u admin -p1234 logistica_global -e "select coste_eur from mantenimientos_flota limit 1;"
\end{lstlisting}

Comprobamos que hay almacenes duplicados:
\begin{lstlisting}[language=bash]
mysql -u admin -p1234 logistica_global -e "select cod_almacen,count(id) from almacenes group by cod_almacen having count(id) > 1 limit 2;"
\end{lstlisting}

Restauramos la incremental del ejercicio 2:
\begin{lstlisting}[language=bash]
mysql -u admin -p1234 logistica_global < 02_incremental_ejercicio2.sql
\end{lstlisting}

Comprobamos que NO hay almacenes duplicados:
\begin{lstlisting}[language=bash]
mysql -u admin -p1234 logistica_global -e "select cod_almacen,count(id) from almacenes group by cod_almacen having count(id) > 1 limit 2;"
\end{lstlisting}

Para probar la diferencial, restauramos la base de datos inicial:
\begin{lstlisting}[language=bash]
mysql -u admin -p1234 logistica_global < 00_copia_inicial.sql
mysql -u admin -p1234 logistica_global < 03_diferencial_ejers_1_y_2.sql 
\end{lstlisting}

Finalmente, restauramos la integridad del ejercicio 3:
\begin{lstlisting}[language=bash]
mysql -u admin -p1234 logistica_global < 03_incremental_ejercicio3.sql 
\end{lstlisting}

Comprobamos que no hay clientes inválidos:
\begin{lstlisting}[language=bash]
mysql -u admin -p1234 logistica_global -e "select count(*) from envios where cliente_id not in (select id from clientes);"
\end{lstlisting}

\end{document}
